Investment has been seen as an important part of explaining the business cycle for a long time. It remains the most volatile variable in the macroeconomy. Over the past few decades, the role of investment in the business cycle has been continuously discussed, with new model environments being proposed \citep{greenwoodInvestmentCapacityUtilization1988,greenwoodLongRunImplicationsInvestmentSpecific1997a,fisherDynamicEffectsNeutral2006,justinianoInvestmentShocksBusiness2010,winberryLumpyInvestmentBusiness2021}.

The discussion of investment in the business cycle is mainly based on investment shocks on the supply side \cite{fisherDynamicEffectsNeutral2006, justinianoInvestmentShocksBusiness2010}, such as the investment specific technology (IST) shock. However, with the development of modeling techniques in heterogeneous models, an investment demand-side shock is also proposed by \cite{nireiOriginsPropagationAnimal}. Unlike supply-side shocks, the demand-side shock affects the firm's investment decision. Moreover, under certain calibrations, it has been shown that the investment demand shock can produce comovement between consumption and inflation, which is otherwise hard to interpret using supply-side shocks. Although this comovement effect is preferred to explain the data, this shock still lacks empirical evidence to support it. Following \cite{justinianoInvestmentShocksBusiness2010}, this paper uses an estimated DSGE model to quantify the influence of the investment demand shock in the business cycle.

As the origin of the investment demand shock stems back to the lumpy investment assumption and the idiosyncratic productivity shock, which works only in models with heterogeneous firms, the shock should be simplified in a representative-agent model. To put this shock in a representative-agent model, we use a two-step solution to push up investment demand.

Using the two-step solution method, I build a DSGE model with both an IST shock and an investment demand shock and estimate it using Bayesian methods. The results show that investment demand shock can be another competitive shock driving the business cycle. Investment demand shock alone contributes about 40\% of the forecast error variance of output, and about 50\% of investment.

To further address concerns about identification problems and oversimplification of the model specification, I also use relative price data to further identify the IST shock, and then use a more standard model similar to \cite{smetsShocksFrictionsUS2007} to solidify the results estimated from the baseline model. These results provide empirical support to quantify the influence of the investment demand shock, and to further investigate its role in the business cycle.